% ==========================================================================
% Conclusions
% ==========================================================================

\section{Conclusions}

In this work, we investigated whether physics-informed machine learning
can compensate for the loss of spatial information introduced during
aggressive reduction of radiotherapy PDE models to low-dimensional ODE
surrogates.

The experiments showed that classical reduced ODE models absorb
unresolved spatial effects into distorted effective parameters, leading
to degraded long-term prediction accuracy. Purely data-driven NODE
models improved trajectory fitting during assimilation but exhibited
unstable extrapolation outside the training interval.
In contrast, the proposed PI-NODE framework combining mechanistic
radiotherapy dynamics with a learnable residual closure term achieved
substantially improved predictive stability across all investigated
beam--tumor configurations. From the Scientific Machine Learning
perspective, the resulting formulation may be interpreted as a
reduced-order UDE-based closure model in which the residual neural
component partially compensates for unresolved spatial damage dynamics
eliminated during PDE-to-ODE reduction.
An additional important result concerns computational efficiency. All
surrogate models reduced prediction costs by several orders of magnitude
compared to the full PDE simulations, making such approaches promising
for future interactive therapy-planning and AI-assisted treatment-design
frameworks.

\subsection*{What the Real-Cohort Evaluation Adds}

Repeating the study on 117 real glioma patients
(Section~\ref{sec:realdata}) confirms the central claim, sharpens it, and
qualifies one of its assumptions.

\paragraph{Closure learning wins where the paper says it should --- by a factor
of two to three.} On the two geometry-mismatched configurations, a learned
closure reaches 9.8\,\% and 9.1\,\% forecast
error against 25.9\,\% and 18.9\,\% for the
best purely mechanistic surrogate. The effect the closure was introduced to
capture is real, is larger on patient anatomy than on the two-dimensional
phantom, and is not explained away by any of the corrections below.

\paragraph{But a single physics--closure balance is not enough.} The optimum
inverts with the irradiation geometry: on the covered configurations the purely
mechanistic model wins and an unconstrained closure fails outright, and on the
mismatched ones the ordering reverses --- the realised physics share of the
per-configuration winner runs from $\Phi\approx1$ to $\Phi\approx0$ across the
same cohort. No constant $(\omega,s_{r})$, and no constant $\lambda$, can be
optimal for both. This is why the useful object is the state-dependent gate of
Eq.~\eqref{eq:gate} rather than a weight to be tuned once, and it is a
conclusion a weight ablation on a single configuration cannot reach.

\paragraph{The weights of Eq.~\eqref{eq:pinode_weighted} are not a mixing
coefficient.} $\omega$ and $s_{r}$ are exact gauge freedoms: each can be
absorbed, to machine precision, into the parameters of its own branch. They
cannot be normalised to sum to one because both branches sweep out cones in
function space, so there is no normalisation to impose. The practical
consequence is that an $(\omega,s_{r})$ ablation reports the optimisation path
rather than the strength of the physics, and that the fitted $\rho,\gamma,\mu$
are gauge coordinates rather than rates. Replacing the pair by a bounded
closure, an anchored backbone and an orthogonality constraint
(Eq.~\eqref{eq:cpinode}) yields a single coefficient whose value matches the
balance the fitted model actually adopts to within a few per cent, and for which
classical linear interpolation between the mechanistic and learned models is
meaningful, over the full $0$--$1$ range rather than the
0.41--0.90 that $(\omega,s_{r})$ can actually reach.

The cost of that identifiability is measured rather than assumed, and the
ablation localises it. The anchor and the orthogonality projection are not the
expensive parts --- within the gated family they \emph{improve} accuracy
(Section~\ref{sec:results_gate}). What costs is the $\tanh$ bound, the one
ingredient that removes the $s_{r}$ gauge: the best convex setting reaches
20.19\,\% against 14.47\,\% for the best cell of the original
grid, and the strongest models in this study are ones with an \emph{unbounded}
closure supplying $\sim90\,\%$ of the vector field.

\paragraph{The reduced growth law is the cheapest available improvement.} A
constant-speed invasion front implies $\dot y\propto y^{2/3}$
(Eq.~\eqref{eq:volumetric_rate}), and the reference trajectories confirm it
(median exponent $0.81$). Adopting that exponent improves the purely
mechanistic surrogate from 25.62\,\% to 23.17\,\% with no additional
parameters, no network and no additional training cost --- and it improves the
PI-NODE's reliability as well, cutting its inter-quartile spread by nearly half.
It is worth being precise about the mechanism: the oracle residuals of the two
backbones are nearly identical (6.96\,\% and 6.79\,\%), so the
volumetric law does not let the model \emph{represent} more --- it makes
extrapolation from a short window better conditioned.

\paragraph{Closure learning is bounded by estimation, not by expressiveness.}
Fitted to the entire noise-free trajectory, the reduced form already reaches
$\approx6.79$\,\% error; the assimilated models are several times
worse. The gap is therefore dominated by what twenty noisy samples cannot
determine. This explains a result that runs against the usual expectation: a
\emph{small} unconstrained residual correction is the most damaging
configuration, because it is flexible enough to absorb in-window noise and its
error then compounds across the forecast. Flexibility helps only when it is
constrained --- bounded in magnitude, anchored against reabsorption into the
mechanistic parameters, and projected off their sensitivities.

\paragraph{The balance should be set, not learned.} A state-dependent gate is
the natural way to exploit the regime dependence above, and it does not work:
the freer the gate, the worse the forecast, with a completely unrestrained gate
settling at a uniform intermediate balance and producing the worst result of its
ablation. Nor is this a missing-input problem --- giving the gate the exact
scalar that separates the regimes changes the forecast by at most $0.24$
percentage points (Section~\ref{sec:results_gate}). The obstruction is the
objective: every setting of the gate fits the assimilation window to the
observation-noise floor, so the regime distinction has no in-window signature at
all, and the payoff lies entirely in a forecast window the loss never sees. No
parameterisation and no side information can recover a schedule that the
training signal does not contain. Since the treatment geometry is known at
planning time and the right choice is worth a factor of two to three, the
practical recommendation is to \emph{set} the physics--closure balance from the
plan rather than fit it: physics-dominated for a conformal, well-covering field;
closure authority for an under-sized or displaced one. Learning it instead
requires a forecast-window objective, and therefore a longitudinal cohort with a
post-treatment observation --- a data-collection requirement rather than a
modelling one.

\paragraph{Limitations.} The reference dynamics are simulated, not measured: the
released dataset provides two static timepoints and no dose-over-time signal, so
a dense response curve cannot be observed. Patient anatomy, tumour geometry and
invasiveness are imaging-derived, but the absolute growth timescale,
radiosensitivity and hypoxic radioresistance are assigned per patient and are
recoverable only through assimilation. What is established is therefore a
comparison of reduced-model \emph{forms} on patient-specific geometry, not a
patient-calibrated clinical prediction, and the recurrence segmentations
available in the dataset remain an untouched validation target.
Future work will focus on more realistic multiscale tumor models
(e.g., similar to \cite{Klusek2019TumorGPU}), uncertainty-aware
surrogate dynamics, explicit latent damage modeling, and integration
with clinical knowledge systems and LLM-driven therapy-design
environments.
